% !PRINTED 397
\weberpage{397}
That this number is really the exact number of the genera we shall prove later. We further remark on this the following:

The fact that for every $\Delta$ there are numbers $a$ for which $(\Delta, a) = -1$ follows from the law of reciprocity for quadratic residues. If one can establish the existence of such numbers $a$ without presupposing this law, then conversely a proof of the law of reciprocity can be derived from it. In this way \gesperrt{Gauss} derived his second proof of this law. (Cf. \gesperrt{Dirichlet-Dedekind}, \S\,152\,ff.)

If $r$ is a norm residue and $\delta$ any stem divisor of $D$, then
\weq{6}{N(\omega) = \frac{x^2 - y^2 D}{4} \equiv r \ (\mathrm{mod}\ \delta)}
and consequently, if $r$ is relatively prime to $\delta$,
\weq{7}{(\delta, r) = +1.}

For odd $\delta$ this is evident; for the even ones it follows easily in the various cases of \S\,108. We therefore set quite generally, when $z$ is relatively prime to $D$:
\weq{8}{\chi(z) = (\delta, z).}

Then $(\delta, r) = 1$ for every number $r$ of the group of absolute norm residues, and in general
\weq{9}{\chi(z) \chi(z_1) = \chi(z z_1).}

For all ideals $\mathfrak{a}$ of the \gesperrt{principal genus} which are coprime to $D$, we have
\[
\chi[N(\mathfrak{a})] = +1,
\]
and from this it follows that $\chi[N(\mathfrak{a})]$ has \gesperrt{one and the same value $\chi(G)$ for all ideals $\mathfrak{a}$ of one genus}.

\gesperrt{These functions are therefore the characters of the group of the genera.}

\begin{center}\rule{0.32\linewidth}{0.4pt}\end{center}
